\chapter{Fundamentos teóricos}\label{ch:theoretical_background}
%\chapter{Theoretical Background}\label{ch:theoretical_background}

%%%%%%%% SECTION 1
%%$$$$$$
\section{Ciclones extratropicales en el Atlántico Sur}\label{sec:tb_fundamentos}
%\section{Extratropical Cyclones in the South Atlantic Basin}\label{sec:tb_fundamentos}
% La teoría clásica define a los ciclones extratropicales como vórtices de núcleo frío impulsados por inestabilidad baroclínica \citep{hoskins2005new}, la caracterización de los sistemas en el Atlántico Sur requiere matices adicionales. 
% Si bien la teoría clásica (e.g., modelo noruego) define a los ciclones extratropicales como vórtices de núcleo frío, y climatologías robustas como la de \citet{hoskins2005new} demuestran que su génesis global está dominada por la fuerte inestabilidad baroclínica de latitudes medias, la caracterización de estos sistemas en la región específica del Atlántico Sur requiere matices dinámicos adicionales.

La teoría clásica (e.g., modelo noruego) define a los ciclones extratropicales como vórtices de núcleo frío, y climatologías robustas como la de \citet{hoskins2005new} demuestran que su génesis global está dominada por la fuerte inestabilidad baroclínica de latitudes medias, la caracterización de estos sistemas en la región específica del Atlántico Sur requiere matices dinámicos adicionales. En este contexto, la transición de una perspectiva euleriana a una lagrangiana resulta crítica; \citet{sinclair1994objective} demostró que el uso de la vorticidad relativa a 850 hPa es físicamente superior a la presión a nivel del mar para identificar sistemas en etapas tempranas, ya que permite aislar la rotación de mesoescala frente al intenso flujo zonal del Hemisferio Sur.
%%%%
%%%
%%
%
En su tesis doctoral, \citet{coutodesouza2024thesis} realiza una revisión exhaustiva de la física de estos sistemas, concluyendo que la definición adiabática estándar es insuficiente para la región.

% TEXTO LITERAL GOZZO & ROCHA (2013) p. 917: "The simulations indicate that in the presence of these fluxes [latent and sensible] the cyclone derwent greater intensification, had a longer life time and longer trajectory, and presented a typical southeastward movement."

% Si bien la inestabilidad baroclínica proporciona el entorno favorable para la génesis, la evolución y la intensidad del sistema son moduladas por procesos diabáticos asociados a la liberación de calor latente \citep{reboita2022from}. En este sentido, \citet{gozzo2013air} demuestran que los flujos de calor en la interfaz aire-mar no solo profundizan el centro de baja presión, sino que son determinantes para el desarrollo de estructuras de mesoescala específicas, como la seclusión cálida. Por tanto, el análisis de los extremos de superficie debe considerar que el acoplamiento océano-atmósfera altera la distribución de la precipitación a lo largo del frente cálido, modificando la simetría térmica y dinámica del vórtice durante su ciclo de vida.

Si bien la inestabilidad baroclínica proporciona el entorno favorable para la génesis, la evolución y la intensidad del sistema puede ser modulada por procesos diabáticos asociados a la liberación de calor latente, especialmente en eventos de transición o con características híbridas. Este comportamiento es ilustrado por los experimentos de sensibilidad numérica presentados por \citet{reboita2022from}, donde la supresión de los flujos de calor latente altera la estructura térmica y la intensidad del vórtice.
% TEXTO LITERAL HART (2003) p. 585: "A cyclone's life cycle can be analyzed within this phase space, providing substantial insight into the cyclone structural evolution."
% Y p. 586: "To broaden the basic binary classification scheme... a continuum Phase Space is proposed."
% Si bien la inestabilidad baroclínica proporciona el entorno favorable para la génesis, la evolución y la intensidad del sistema son moduladas por procesos diabáticos asociados a la liberación de calor latente \citep{reboita2022from}. En este sentido, \citet{gozzo2013air} demuestran que los flujos de calor en la interfaz aire-mar no solo profundizan el centro de baja presión, sino que son determinantes para el desarrollo de estructuras de mesoescala específicas, como la seclusión cálida, y para definir la trayectoria del sistema hacia el sudeste. La existencia de estas diversas configuraciones evidencia que la evolución ciclónica no es un proceso rígido. Tal como formalizó \citet{hart2003cyclone}.


%%%%%%%% SECTION 2
%%
\section{Regiones de ciclogénesis}\label{sec:tb_regiones}

% % TEXTO LITERAL REBOITA (2010) p. 1331: "The development of these systems is associated with the passage of synoptic waves and the baroclinic instability, which is intensified by the strong thermal gradient between the South American continent and the South Atlantic Ocean."
A partir del análisis de la vorticidad relativa en niveles bajos, \citet{reboita2010south} establecieron la zonificación climática de estos procesos, delimitando las regiones de la costa de Uruguay y el sur de Brasil, así como la región costera de Argentina, como los principales focos de actividad. Mientras que los sistemas extratropicales puros dominan las latitudes altas, \citet{gozzo2014subtropical} demuestran que en las regiones de latitudes medias y subtropicales (20°S--35°S) existe una alta frecuencia de génesis de ciclones híbridos. La interacción entre el flujo de los oestes, la topografía de los Andes y la Temperatura Superficial del Mar (TSM) define tres regiones de génesis con características físicas diferenciadas \citep{gramcianinov2019properties}: 
% Esta caracterización regional es fundamental para la presente investigación, ya que permite entender cómo la ubicación de la génesis condiciona las trayectorias resultantes y, en consecuencia, la asimetría espacial de los extremos de viento y precipitación que se manifiestan en el continente.

% TEXTO LITERAL GOZZO ET AL. (2014) p. 2242: "The geographical distribution of SC genesis shows that most of the systems form over the ocean near the Brazilian coast, between 20 and 35S." Y p. 2239: "SCs in the South Atlantic are more frequent in the summer and autumn."
%
\subsection{Sur-sureste de Brasil (SBR)}\label{subsec:tb_sbr}
%\subsection{Southern-Southeastern Brazil (SBR)}\label{subsec:tb_sbr}
Según \citet{reboita2022from}, RG1 (análoga a la SBR) constituye una de las regiones preferenciales de ciclogénesis en la costa sur/sudeste de Brasil, y el caso de estudio del ciclón Raoni ilustra cómo los sistemas que actúan en este sector pueden presentar características subtropicales durante su ciclo de vida. Son conocidos asi estos ciclones como "híbridos", porque aunque se inician por inestabilidad baroclínica, su intensificación está modulada por procesos diabáticos y la inestabilidad de la capa límite marina, lo que favorece el desarrollo de convección en el sector cálido \citep{marrafon2022classificacao}. 
El análisis energético de \citet{coutodesouza2024thesis} revela que, a diferencia de las regiones australes, los ciclones en SBR presentan una conversión de energía "mixta", donde los procesos diabáticos (liberación de calor latente) juegan un rol tan importante como la baroclinicidad.
Proyecciones con modelos climáticos regionales de alta resolución confirman este comportamiento: \citet{gramcianinov2024early} demuestran que los RCMs aportan valor añadido respecto a los GCMs precisamente en la representación de las estructuras de mesoscala de advección cálida y convergencia de humedad en capas bajas durante la ciclogénesis temprana en SBR, mecanismos que los modelos de baja resolución sistemáticamente suavizan. Este resultado valida que la dinámica híbrida de SBR no es un artefacto estadístico sino una señal física robusta, detectable incluso en marcos de referencia lagrangianos centrados en el sistema.

% Proyecciones con modelos climáticos regionales de alta resolución confirman este comportamiento: \citet{gramcianinov2024early} demuestran que los RCMs aportan valor añadido respecto a los GCMs precisamente en la representación de las estructuras de mesoscala de advección cálida y convergencia de humedad en capas bajas durante la ciclogénesis temprana en SBR, mecanismos que los modelos de baja resolución sistemáticamente suavizan. Este resultado valida que la dinámica híbrida de SBR no es un artefacto estadístico sino una señal física robusta, detectable incluso en marcos de referencia lagrangianos centrados en el sistema.

\subsection{Cuenca del Río de la Plata (LPB)}\label{subsec:tb_lpb}
%\subsection{La Plata Basin(LPB)}\label{subsec:tb_lpb}
La región comprendida entre $30^\circ$S y $40^\circ$S constituye uno de los principales \textit{hotspots} de ciclogénesis en el Atlántico Sudoccidental. Su mecanismo dominante es la ciclogénesis de sotavento: como establecen \citet{gan1994influence}, el flujo de los oestes experimenta una compresión y posterior estiramiento vertical de la columna de atmosfera al atravesar los Andes. Por conservación de la vorticidad potencial, este estiramiento en el flanco oriental induce una caída de presión en superficie que favorece la iniciación de la rotación ciclónica. 

En este contexto, \citet{reboita2010south} muestran que esta región se caracteriza por una alta frecuencia de sistemas ciclónicos y una interacción con el transporte de humedad desde latitudes tropicales, mientras que el aporte de aire cálido y húmedo asociado al Chorro de Capas Bajas (SALLJ) resulta determinante para la intensidad de la precipitación \citep{crespo2020potential}.

% Consistentemente, \citet{gramcianinov2024early} confirman mediante modelos regionales que la convergencia de humedad en capas bajas constituye el mecanismo dominante de intensificación en LPB, proyectando además que bajo escenarios de calentamiento global este forzante se reforzará hacia fines de siglo. Si bien la frecuencia total de ciclones en la región disminuirá, los sistemas que persistan exhibirán flujos de humedad y advección de calor más intensos en sus etapas tempranas, lo que sugiere una mayor severidad potencial de los eventos de precipitación extrema asociados.
Consistentemente, \citet{gramcianinov2024early} confirman mediante modelos regionales que la convergencia de humedad en capas bajas constituye el mecanismo dominante de intensificación en LPB, proyectando además que bajo escenarios de calentamiento global este forzante se reforzará hacia fines de siglo. Si bien la frecuencia total de ciclones en la región disminuirá, los sistemas que persistan exhibirán flujos de humedad y advección de calor más intensos en sus etapas tempranas, lo que sugiere una mayor severidad potencial de los eventos de precipitación extrema asociados.

\subsection{Argentina (ARG)}\label{subsec:tb_arg}
El sector austral ($>40^\circ$S, o su similar RG3) responde a un régimen dinámico clásico dominado por la baroclinicidad de gran escala. 

% A diferencia de los focos subtropicales, \citet{hoskins2005new} describen esta zona como gobernada por la inestabilidad baroclínica, donde la génesis depende del chorro polar y del paso de ondas de Rossby. 
A diferencia de los focos subtropicales, \citet{hoskins2005new} demuestran que el desarrollo en las latitudes australes está estrictamente acoplado en toda la columna troposférica, siendo gobernado por la profunda inestabilidad baroclínica asociada al chorro polar y la propagación de ondas de Rossby a lo largo del storm track de alta latitud.

Aquí, \citet{simmonds2000variability} indican que los sistemas ciclónicos dependen menos de forzantes locales de superficie y más de la dinámica de la alta troposfera. Como resultado, estos sistemas tienden a presentar una evolución más rápida y directamente acoplada al flujo de niveles altos, en contraste con las regiones subtropicales del sudeste de Brasil \citep{vera2002cold}.

%%%%%%%% SECTION 3
%%
\section{Ciclo de vida y Modelo conceptual}\label{sec:tb_evolucion}
%\section{Lifecycle and Conceptual Model}\label{sec:tb_evolucion}

\subsection{Modelo conceptual}\label{subsec:tb_modelos}
%\subsection{Conceptual Model}\label{subsec:tb_modelos}
La distribución espacial de las variables en superficie (viento y precipitación) está condicionada por el modelo conceptual que describe la evolución del ciclón. Frente al modelo Noruego clásico, que presenta una oclusión simétrica, los ciclones oceánicos intensos del Atlántico Sur suelen seguir el modelo de Fractura Frontal propuesto por \citet{shapiro1990life}. En este paradigma, el frente frío se fractura y se separa del centro, permitiendo el aislamiento de una masa de aire cálido en el núcleo (\textit{warm seclusion}).

Para comprender la dinámica interna que rige la asimetría de estas variables, resulta fundamental complementar este paradigma con el modelo de cinturones de transporte \citep{browning1986conceptual}. En su adaptación conceptual para el Hemisferio Sur, Silva Dias y Hallak (2012) detallan cómo el Cinturón Transportador Cálido (\textit{Warm Conveyor Belt}) y su respectiva corriente en chorro de bajo nivel se organizan en torno al vórtice. El ascenso y giro anticiclónico de este flujo cálido modula las estructuras de mesoescala, concentrando las bandas de precipitación y el fuerte cizallamiento del viento (asociado a las ráfagas) por delante del frente frío. Esta configuración cinemática determina que los máximos de viento y precipitación se desacoplen en relacion al centro, un rasgo recurrente documentado en la climatología regional por \citet{reboita2010south}.
%%%

\subsection{Fases del ciclo de vida}\label{subsec:tb_fases}
Para caracterizar la evolución temporal del viento y la precipitación, este estudio adopta la segmentación del ciclo de vida en cuatro etapas dinámicas discretas. Esta clasificación sigue la metodología objetiva del paquete computacional \textit{CycloPhaser} \citep{deSouza2025CycloPhaser}, aplicado por \citet{coutodesouza2024new}, el cual utiliza la tasa de cambio de la vorticidad relativa para definir los siguientes regímenes estructurales. Cabe notar que cada fase no es homogénea internamente: los instantes próximos a sus transiciones comparten rasgos de la etapa adyacente, por lo que la condición estructural más representativa de cada régimen dinámico corresponde al núcleo temporal de la fase, alejado de ambos bordes de transición.

\begin{enumerate}
    \item Incipiente: Corresponde a la aparición inicial del núcleo de vorticidad organizada. En el contexto regional, esta fase está frecuentemente asociada a la interacción entre una vaguada en altura y forzantes de superficie, como la orografía de los Andes o gradientes térmicos costeros.    
    \item Intensificación: Período caracterizado por el rápido crecimiento de la vorticidad ciclónica. Según el análisis energético de \citet{coutodesouza2024thesis}, esta etapa es impulsada por una fuerte retroalimentación entre la inestabilidad dinámica y los flujos de calor latente y sensible desde el océano.    
    \item Madurez: Etapa donde el sistema alcanza su intensidad máxima (pico de vorticidad) y su mayor extensión horizontal. Es durante esta fase donde los ciclones oceánicos intensos consolidan la evolución estructural documentada observacionalmente por \citet{shapiro1999bridge}. En esta configuración, la fuerte deformación del flujo genera la fractura del frente frío y envuelve una masa de aire para que quede aislada en el núcleo del vórtice (\textit{warm seclusion}). Al completarse este proceso las huellas de precipitación y viento quedan desacopladas.  
    \item Decaimiento: Fase final marcada por el debilitamiento progresivo del vórtice. Este proceso ocurre debido a la fricción superficial y al cese de la conversión de energía baroclínica una vez que el sistema se ha ocluido completamente o ha entrado en una zona barotrópica.
\end{enumerate}


%%%%%%%% SECTION 4
%%
\section{Vorticidad relativa}\label{sec:tb_vorticidad}
%\section{Relative Vorticity}\label{sec:tb_vorticidad}
Debido a la extensión oceánica del Hemisferio Sur, los ciclones extratropicales se desarrollan inmersos en un flujo de los oestes. Esta fuerte corriente de fondo hace que los sistemas se trasladen con gran rapidez y que, en sus etapas iniciales, la caída de presión se vea enmascarada, presentándose a menudo como vaguadas abiertas en lugar de centros cerrados \citep{sinclair1994objective}. Dada esta rapidez de traslación y complejidad térmica, la presión al nivel del mar (MSLP) resulta una métrica inadecuada para la detección temprana. 

Para resolver este problema y aislar la circulación del sistema, la recomendación metodológica de \citet{sinclair1997objective} es utilizar la vorticidad relativa en 850 hPa ($\zeta_{850}$) como variable de seguimiento. Físicamente, la componente vertical de la vorticidad relativa se define como:

\begin{equation}
\zeta = \mathbf{k} \cdot (\nabla \times \mathbf{V}) = \frac{\partial v}{\partial x} - \frac{\partial u}{\partial y}
\end{equation}
% TEXTO LITERAL CORNÉR ET AL. (2025) p. 219: "The sparse PCA results allow for an intuitive classification of ETCs based on their dominant intensity characteristics. [...] In contrast to the traditional single-measure approach, our results highlight that intensity is multi-faceted." Y p. 209 (Tabla 1): Incluye "Relative vorticity at 850 hPa" como una de las variables clave analizadas.
%
Donde $u$ y $v$ son las componentes zonal y meridional del viento, respectivamente. Según \citet{hoskins2005new}, el uso de $\zeta$ permite filtrar el flujo de gran escala y capturar la rotación inducida por estos forzantes de mesoescala antes de que se manifiesten en el campo de presión, permitiendo una identificación más precisa de las fases incipientes, ademas la literatura reciente, como en \citet{gramcianinov2019properties} utiliza $\zeta_{850}$ para identificar con precisión las regiones de ciclogénesis. Asimismo, \citet{padilhareinke2024objective} validan que los algoritmos basados en este nivel isobárico son más eficientes para capturar la posición del ciclón durante su etapa de mayor intensidad. 

Recientemente, el uso de $\zeta_{850}$ se ha extendido más allá del simple seguimiento (tracking) hacia la caracterización estructural. \citet{coutodesouza2024new} demuestran que la evolución temporal de esta variable es el predictor más confiable para segmentar objetivamente el ciclo de vida en fases dinámicas (incipiente, intensificación, madurez y decaimiento). En un contexto global el análisis de clústeres de \citet{corner2025classification} ratifican a la vorticidad relativa en 850 hPa como una de las métricas irreductibles necesarias para describir la intensidad y la estructura de los ciclones extratropicales en el hemisferio norte.

El uso de umbrales percentílicos sobre la distribución de $\zeta_{850}$ para estratificar ciclones por intensidad dinámica cuenta con respaldo metodológico sólido. \citet{priestley2022improved} aplican explícitamente este criterio seleccionando el 10\% de ciclones con mayor vorticidad ciclónica en su pico de intensidad para construir compositos representativos de la estructura ciclónica en el Hemisferio Sur, demostrando que esta selección permite aislar los patrones espaciales más robustos y físicamente interpretables. En el contexto regional del Atlántico Sudoccidental, \citet{andrade2024composite} documentan de forma complementaria que los sistemas ubicados en el percentil 10 de presión mínima —equivalente dinámico al p90 de vorticidad negativa— exhiben patrones sinópticos diferenciados y una mayor coherencia entre los extremos de viento en superficie y la profundidad del vórtice, validando estadísticamente esta población como una submuestra física y estructuralmente homogénea.


%%%%%%%% SECTION 5
%%
\section{Estructura superficial y extremos}\label{sec:tb_compuestos}
%\section{Surface Structure and Extrems}\label{sec:tb_compuestos}
La complejidad inherente a los ciclones extratropicales del Atlántico Sur requiere superar el análisis univariado tradicional. En esta sección se fundamenta la adopción del paradigma de \textit{evento compuesto} y se describen las herramientas estadísticas avanzadas (MEOF y KDE) necesarias para capturar su estructura asimétrica y la distribución de sus máximos.

\subsection{Asimetría espacial de los Extremos (KDE)}\label{subsec:tb_kde_asimetria}
%\subsection{Spatial Asymmetry of Extrems (KDE)}\label{subsec:tb_kde_asimetria}
Los ciclones extratropicales son sistemas altamente asimétricos, donde los máximos raramente coinciden con el centro de vorticidad. Es importante notar que gran parte de los modelos conceptuales clásicos fueron desarrollados para el Hemisferio Norte; por lo tanto, al analizar el Atlántico Sur, la distribución espacial debe entenderse como una imagen invertida latitudinalmente tomando al ecuador como eje.

\begin{itemize}
%
    \item Precipitación: Está fundamentalmente controlada por el \textit{Warm Conveyor Belt} (WCB). En la literatura clásica del Hemisferio Norte, este flujo se describe ascendiendo típicamente hacia el noreste \citep{blanchard2021warm}. Sin embargo, al trasladar este concepto al Hemisferio Sur, el efecto del espejo hace que el WCB concentre la humedad en el sector sureste del sistema durante su etapa de desarrollo. De acuerdo con el modelo conceptual adaptado para el Hemisferio Sur por Silva Dias y Hallak (2012) a partir de \citet{browning1986conceptual}, el ascenso inclinado de este flujo cálido organiza la precipitación en estructuras de mesoescala bien definidas, generando bandas anchas de lluvia estratiforme por delante del sistema y bandas estrechas de convección vigorosa en la interfaz del frente frío. Además, como demuestran \citet{oertel2021warm}, el WCB no solo produce esta precipitación a gran escala, sino que frecuentemente alberga núcleos de convección intensa embebida. La marcada asimetría de este campo es corroborada empíricamente por \citet{naud2020evaluation}, quienes demuestran que la precipitación en los ciclones oceánicos se concentra abrumadoramente en el sector oriental (ascendente y húmedo), mientras que el flanco occidental permanece subsidente y seco.
%
    \item Precipitación: Está fundamentalmente controlada por el \textit{Warm Conveyor Belt} (WCB). Mientras que en el Hemisferio Norte este flujo asciende hacia el noreste, en el Hemisferio Sur se dirige hacia latitudes más altas, concentrando la humedad en el sector sureste del sistema durante su etapa de desarrollo \citep{blanchard2021warm}. Además, como demuestran \citet{oertel2021warm}, el WCB no solo produce precipitación estratiforme a gran escala, sino que frecuentemente alberga núcleos de convección intensa embebida. La marcada asimetría de este campo es corroborada empíricamente por \citet{naud2020evaluation}, quienes demuestran que la precipitación en los ciclones oceánicos se concentra abrumadoramente en el sector oriental (ascendente y húmedo), mientras que el flanco occidental permanece subsidente y seco.
    \item Viento: Al igual que la precipitación, el campo de viento en superficie presenta una asimetría que cambia con el ciclo de vida del sistema. En el Atlántico Sur, \citet{bitencourt2010relating} y \citet{cardoso2022synoptic} documentan que los vientos máximos tienden a concentrarse en sectores específicos, como el cuadrante noreste. Además, las ráfagas más destructivas suelen estar vinculadas a características de mesoescala cuya posición relativa al centro de vorticidad varía a medida que el ciclón madura \citep{eisenstein2023identification}. Esta heterogeneidad explica por qué la intensidad de los máximos del viento no puede capturarse mediante una única métrica central \citep{corner2025classification}.
\end{itemize}

Esta disociación entre la métrica dinámica central y el impacto en superficie tiene consecuencias directas para el diseño analítico. \citet{andrade2024composite} documentan que, en el Atlántico Sudoccidental, los ciclones más intensos según MSLP presentan trayectorias sistemáticamente desplazadas hacia el sur respecto a los más intensos según viento máximo a 10 m, y que solo el 56\% de los sistemas que superan ambos umbrales (p10 de presión y p90 de viento) simultáneamente afectan la zona costera. Este resultado evidencia que el análisis de los extremos de viento requiere un marco de referencia centrado en el sistema —como el que proporciona el KDE bidimensional— y no un enfoque euleriano fijo, ya que la asimetría espacial de los máximos varía entre sistemas de igual intensidad dinámica central.

Debido a que los máximos de estas variables ocurren desacoplados del centro del ciclon, intentar representarlos mediante promedios simples o grillas fijas tiende a suavizar y ocultar estas estructuras. Para capturar la distribución real de los extremos sin perder información espacial de los ciclones, autores como \citet{priestley2022improved} y \citet{han2025system} proponen el uso de la Estimación de Densidad de Kernel (KDE).

\begin{equation}
\hat{f}(x) = \frac{1}{nh} \sum_{i=1}^{n} K\left(\frac{x - X_i}{h}\right)
\end{equation}

Esta técnica transforma observaciones discretas en una superficie de probabilidad continua, permitiendo mapear la \textit{huella} física del ciclón. Su aplicación es especialmente útil al segmentar el análisis por etapas evolutivas, aprovechando asi los resultados del algoritmo \textit{CycloPhaser} \citep{coutodesouza2024new}, ya que esto puede revelar cómo los campos extremos se expanden y contraen a lo largo del ciclo de vida del sistema.

\subsection{Extremos de viento y precipitación}\label{subsec:tb_compound_def}
%\subsection{Extremes of Wind and Precipitation}\label{subsec:tb_compound_def}
Históricamente, la intensidad de los ciclones se ha caracterizado mediante métricas puntuales, como la presión mínima central o la vorticidad máxima. Sin embargo, la literatura reciente ha adoptado el enfoque de \textit{compound events} (eventos compuestos), definidos formalmente por \citet{zscheischler2018future} como situaciones donde la combinación de múltiples forzantes (p. ej., viento y precipitación extremos) genera magnitudes conjuntas mayores que la suma de sus componentes individuales. En el contexto de la dinámica extratropical, \citet{chen2025characteristics} establecen que la interacción entre el campo de viento y la precipitación no es lineal ni espacialmente homogénea. La ocurrencia conjunta de estos extremos responde a mecanismos físicos acoplados: la intensificación dinámica suele estar modulada por procesos diabáticos (liberación de calor latente), lo que sugiere una dependencia estadística entre ambas variables que varía según la fase de vida del sistema. Por tanto, para caracterizar la dinámica del sistema, se requiere un enfoque multivariado que cuantifique esta covariabilidad física.
No obstante, es necesario reconocer una limitación instrumental en la cuantificación de estos extremos. Evaluaciones sistemáticas de ERA5 sobre variables de superficie asociadas a ciclones demuestran que, si bien el producto presenta un sesgo global reducido para el viento (-0.7\%) y la precipitación (-10.4\%), este sesgo se amplifica significativamente para los eventos más intensos: para el 5\% de ETCs más extremos, la subestimación del viento alcanza el -10.2\% y la de la precipitación el -22.6\% \citep{chen2024evaluation}. Esta tendencia a suavizar los máximos más intensos implica que los resultados obtenidos para el subconjunto p90 deben interpretarse como estimaciones conservadoras de la severidad real de los sistemas más profundos.

\subsection{Descomposición modal de campos (EOF/MEOF)}\label{subsec:tb_math_meof}
%\subsection{Modal Decomposition of Coupled Fields (EOF/MEOF)}\label{subsec:tb_math_meof}
Para aislar patrones coherentes dentro de este acoplamiento físico, el análisis de Funciones Ortogonales Empíricas (EOF) y su extensión multivariada (MEOF) se han establecido como el estándar metodológico. Según \citet{hannachi2007empirical}, esta técnica permite descomponer la variabilidad de un campo en modos ortogonales de varianza máxima, separando la señal física robusta del ruido de mesoescala. Específicamente para eventos compuestos, \citet{liang2018multivariate} demostraron la eficacia del MEOF al aplicarlo sobre un vector de estado normalizado $\mathbf{Z}$, que combina variables con unidades dispares: 
\begin{equation}
\mathbf{Z}(t) = \left[ \frac{X'*{prep}(t)}{\sigma*{prep}}, \frac{X'*{viento}(t)}{\sigma*{viento}} \right]
\end{equation}
Esta normalización permite identificar si los máximos de precipitación y viento covarían espacialmente o si presentan desfases estructurales sistemáticos. Asimismo, enfoques recientes de clasificación objetiva, como el PCA disperso utilizado por \citet{corner2025classification}, validan el uso de espacios de fase reducidos para distinguir entre regímenes dinámicos. 
Este diseño sigue la práctica establecida por estudios de estructura ciclónica como el de \citet{priestley2022improved}, quienes aplican compositing centrado en el sistema al 10\% de ciclones más intensos (en términos de vorticidad) para aislar los patrones estructurales representativos del Hemisferio Sur.
